\documentclass[../../../include/open-logic-section]{subfiles}

\begin{document}

\olfileid{sth}{replacement}{strength}
\olsection{প্রতিস্থাপনের শক্তি}

প্রতিস্থাপনের শক্তি সম্পর্কে একটি সরল পর্যবেক্ষণ
দিয়ে শুরু করি: $\Z$-এর বাইরে না গেলে
$\omega + \omega$-র সমান বা বড় কোনো
ফন নয়মান ক্রমসংখ্যার অস্তিত্ব প্রমাণ করা যায় না।

কেন, তার একটি রূপরেখা দিই।
~$\ZF$-এ কাজ করে $V_{\omega+\omega}$
সেটটি বিবেচনা করি। সেটটি $\Z$-এর একটি
\emph{মডেলের} ভিত্তিসেট হিসেবে কাজ করে।
এটি দেখতে সূত্রের \emph{আপেক্ষিকীকরণ}-এর জন্য
কিছু সংকেত নিই:
\begin{defn}\ollabel{formularelativization}
	যেকোনো সেট $M$ ও সূত্র $\phi$-এর জন্য
	$\phi^M$ হলো $\phi$-র সব পরিমাণায়ককে
	$M$-এ সীমাবদ্ধ করে পাওয়া সূত্র। অর্থাৎ
	``$\lexists[x]$''-এর জায়গায়
	``$(\lexists[x \in M])$'' এবং
	``$\lforall[x]$''-এর জায়গায়
	``$(\lforall[x \in M])$'' বসাই।
\end{defn}\noindent
দেখানো যায়, $\Z$-এর প্রত্যেক স্বতঃসিদ্ধ
$\phi$-এর জন্য $\ZF \vdash
\phi^{V_{\omega+\omega}}$। কিন্তু
\olref[spine][rank]{ordsetrankalpha} অনুযায়ী
$\omega+\omega$, $V_{\omega+\omega}$-র
\emph{সদস্য নয়}। তাই $\omega+\omega$-র
অনস্তিত্বের সঙ্গে $\Z$ সঙ্গতিপূর্ণ।

এই কারণেই \olref[ordinals][replacement]{sec}-এ
বলেছিলাম, প্রতিস্থাপন ছাড়া
\olref[ordinals][ordtype]{thmOrdinalRepresentation}
প্রমাণ করা যায় না। কারণ $\Z$-এর মধ্যেই
এমন একটি স্পষ্ট সুক্রমবিন্যাস সংজ্ঞায়িত করা
সহজ, যার ক্রমপ্রকার স্বজ্ঞায় $\omega+\omega$
\emph{হওয়া উচিত}। বস্তুত,
\olref[ordinals][idea]{sec}-এ স্বাভাবিক সংখ্যার
উপর নিচের ক্রমটি দিয়ে এর একটি অনানুষ্ঠানিক
উদাহরণ দিয়েছি:
\begin{align*}
	n \lessdot m \text{ যদি ও কেবল যদি }&\text{হয় }n < m\text{ এবং }m-n\text{ জোড়,}\\
	& \text{অথবা $n$ জোড় এবং $m$ বিজোড়।}
\end{align*}
কিন্তু $\omega+\omega$-র অস্তিত্ব না থাকলে
এই সুক্রমবিন্যাস কোনো ক্রমসংখ্যার সমরূপ হবে
না। সুতরাং $\Z$
\olref[ordinals][ordtype]{thmOrdinalRepresentation}
\emph{প্রমাণ করে না}।

উল্টো দিক থেকে দেখি: প্রতিস্থাপন দিয়ে
$\omega+\omega$-র অস্তিত্ব প্রমাণ করা যায়;
ফলে এটি $V_{\omega+\omega}$-র অস্তিত্বও
প্রমাণ করতে দেয়। এখানেই শেষ নয়।
আমরা সংজ্ঞায়িত করতে পারি এমন
\emph{যেকোনো} সুক্রমবিন্যাসের জন্য
\olref[ordinals][ordtype]{thmOrdinalRepresentation}
বলে, সেটির সমরূপ কোনো $\alpha$ আছে;
ফলে $V_\alpha$-ও আছে। তাই সোজা কথায়,
প্রতিস্থাপন নিশ্চিত করে যে সেটের স্তরক্রম
\emph{অনেক উঁচু} পর্যন্ত বিস্তৃত।

পরের কয়েকটি বিভাগে, এবং আবার
\olref[card-arithmetic][fix]{sec}-এ,
প্রতিস্থাপন স্তরক্রমকে ঠিক \emph{কতটা}
উঁচু হতে বাধ্য করে, তা ভালোভাবে বুঝব।
আপাতত সরল কথাটি এই: প্রতিস্থাপনের
ন্যায্যতা \emph{সত্যিই} প্রতিষ্ঠা করা দরকার!

\end{document}
